% Informations générales du document
\newcommand{\myAuthor}{Étienne Caulier}
\newcommand{\myLesson}{TPI}
\newcommand{\myEmail}{etienne.clr@eduge.ch}
\newcommand{\myTitle}{Apprentissage \LaTeX}
\newcommand{\myClass}{}
\newcommand{\mySchool}{}
\newcommand{\mySubjectStudies}{Informatique}

% Imports des packages
\usepackage[french]{babel}
\usepackage[T1]{fontenc}
\usepackage[utf8]{inputenc}
\usepackage{adjustbox}
\usepackage{array}
\usepackage{booktabs}
\usepackage{fancyhdr, fancyvrb}
\usepackage{float}
\usepackage{forest}
\usepackage{geometry}
\usepackage{graphicx}
\usepackage{helvet}
\usepackage{hyperref}
\usepackage{lastpage}
\usepackage{listings}
\usepackage{longtable}
\usepackage{makecell}
\usepackage{multirow}
\usepackage{pdfpages}
\usepackage{pifont}
\usepackage{titlesec}
\usepackage{xcolor}
\usepackage[table]{xcolor}

% Config packages
\geometry{margin=2.5cm}
\renewcommand{\familydefault}{\sfdefault} % police sans-serif

% Tableau
\setcellgapes{1pt}
\makegapedcells
\newcolumntype{R}[1]{>{\raggedleft\arraybackslash }b{#1}}
\newcolumntype{L}[1]{>{\raggedright\arraybackslash }b{#1}}
\newcolumntype{C}[1]{>{\centering\arraybackslash }b{#1}}
\renewcommand{\arraystretch}{1.4}

% Initialise les en-têtes et pieds de page
\pagestyle{fancy}
\fancyhf{}

% Import custom packages
%======================================================================
% includes/inc_colors.tex
% Palette de couleurs centralisée pour tous les includes
%======================================================================

\usepackage{xcolor}

%----------------------------------------------------------------------
% Couleur de texte "atténué" (muted)
%----------------------------------------------------------------------
\definecolor{mutedColor}{gray}{0.45} % texte secondaire (pied de page, légendes…)

%----------------------------------------------------------------------
% Gris / neutres de base
%----------------------------------------------------------------------
\definecolor{cfptBlack}{rgb}{0.00,0.00,0.00}
\definecolor{cfptDarkGray}{gray}{0.25}
\definecolor{cfptGray}{gray}{0.50}
\definecolor{cfptLightGray}{gray}{0.90}
\definecolor{cfptVeryLightGray}{gray}{0.98}

% Alias "historiques" pour compatibilité
\colorlet{lightgray}{cfptVeryLightGray}
\colorlet{gray}{cfptGray}

%======================================================================
% Palette pour les listings de code
%======================================================================

% Fond des blocs de code
\colorlet{codeBackground}{cfptVeryLightGray}

% Couleurs sémantiques pour la coloration syntaxique
\definecolor{codeKeyword}{rgb}{0.00,0.00,0.80}    % bleu foncé
\definecolor{codeIdentifier}{rgb}{0.00,0.60,0.00} % vert
\definecolor{codeString}{rgb}{0.80,0.00,0.00}     % rouge
\definecolor{codeComment}{rgb}{0.50,0.50,0.50}    % gris
\definecolor{codeAccent}{rgb}{0.58,0.00,0.82}     % violet (true/false/null)
\definecolor{codeOperator}{rgb}{0.90,0.45,0.00}   % orange (opérateurs, ponctuation)
\definecolor{codeType}{rgb}{0.20,0.60,0.60}       % cyan pour types/classes

% Alias "historiques" pour compatibilité avec les anciens includes
\colorlet{dkblue}{codeKeyword}
\colorlet{dkgreen}{codeIdentifier}
\colorlet{red}{codeString}
\colorlet{purple}{codeAccent}
\colorlet{orange}{codeOperator}   % !!! réservé à la coloration de code
\colorlet{class}{codeType}
\colorlet{darkgreen}{codeIdentifier}
\colorlet{darkred}{codeString}

%======================================================================
% Palette pour les boîtes (info / success / warning / danger / etc.)
%  -> Inspirée de Bootstrap (blue / green / amber / red) avec fonds doux.
%======================================================================

% Couleurs pour la boîte par défaut
\definecolor{defaultBackground}{RGB}{245,245,245} 		% gris très clair
\colorlet{defaultColor}{gray}                     		% bordure grise
\colorlet{defaultTitleBackground}{gray!15}        		% fond du titre légèrement grisé

% INFO  : messages neutres / informatifs
\definecolor{infoColor}{rgb}{0.00,0.45,0.80}            % bleu
\definecolor{infoBackground}{rgb}{0.91,0.96,1.00}       % bleu très clair
\definecolor{infoTitleBackground}{rgb}{0.00,0.45,0.80}  % bandeau titre

% SUCCESS : succès, validation, confirmation
\definecolor{successColor}{rgb}{0.17,0.63,0.34}               % vert
\definecolor{successBackground}{rgb}{0.91,0.98,0.92}          % vert très clair
\definecolor{successTitleBackground}{rgb}{0.17,0.63,0.34}

% WARNING : avertissement, point de vigilance
\definecolor{warningColor}{rgb}{0.98,0.76,0.18}               % jaune / amber
\definecolor{warningBackground}{rgb}{1.00,0.97,0.88}          % jaune très clair
\definecolor{warningTitleBackground}{rgb}{0.98,0.76,0.18}

% DANGER : erreur critique, interdiction
\definecolor{dangerColor}{rgb}{0.83,0.25,0.25}                % rouge
\definecolor{dangerBackground}{rgb}{1.00,0.93,0.94}           % rouge très clair
\definecolor{dangerTitleBackground}{rgb}{0.83,0.25,0.25}

% TEACHER : notes internes formateur
\definecolor{teacherColor}{rgb}{0.55,0.27,0.07}         	% brun soutenu
\definecolor{teacherBackground}{rgb}{0.98,0.95,0.90}    	% beige très clair
\definecolor{teacherTitleBackground}{rgb}{0.55,0.27,0.07}

% DOWNLOAD : encadrés de téléchargement / ressources
\definecolor{downloadColor}{rgb}{0.00,0.40,0.80}              % bleu soutenu
\definecolor{downloadBackground}{rgb}{0.92,0.96,1.00}         % bleu clair
\definecolor{downloadTitleBackground}{rgb}{0.00,0.40,0.80}

% AIDE : indices, coups de pouce
\definecolor{aideColor}{rgb}{0.40,0.23,0.72}                  % indigo / violet
\definecolor{aideBackground}{rgb}{0.95,0.93,1.00}             % indigo très clair
\definecolor{aideTitleBackground}{rgb}{0.40,0.23,0.72}

% OBJECTIF : objectifs pédagogiques, buts de séance
\definecolor{objectifColor}{rgb}{1.00,0.58,0.00}              % orange franc
\definecolor{objectifBackground}{rgb}{1.00,0.95,0.90}         % orange très clair
\definecolor{objectifTitleBackground}{rgb}{1.00,0.58,0.00}

% VALIDATION : critères de validation, checklist finale
\definecolor{validationColor}{rgb}{0.00,0.58,0.51}            % vert/bleu (teal)
\definecolor{validationBackground}{rgb}{0.90,0.97,0.96}       % teal très clair
\definecolor{validationTitleBackground}{rgb}{0.00,0.58,0.51}

% (Pas d'alias ici sur 'orange' ou 'cyan' pour éviter de casser les listings)

%======================================================================
% USAGE :
%   \begin{boxDefault}[Titre optionnel] ... \end{boxDefault}   % boîte grise sans icône
%   \begin{boxInfo}[Titre optionnel] ... \end{boxInfo}
%   \begin{boxSuccess}[Bravo !] ... \end{boxSuccess}
%   \begin{boxWarning}[Attention] ... \end{boxWarning}
%   \begin{boxDanger}[Important] ... \end{boxDanger}
%   \begin{boxTeacher}[Note pour l'enseignant] ... \end{boxTeacher}
%   \begin{boxDownload}[Ressources] ... \end{boxDownload}
%   \begin{boxHelp}[Indice] ... \end{boxHelp}
%   \begin{boxObjectif}[Objectifs] ... \end{boxObjectif}
%   \begin{boxValidation}[Critères] ... \end{boxValidation}
%======================================================================

%----------------------------------------------------------------------
% Gestion du chemin des icônes (\DocDisplayBoxPictures)
%----------------------------------------------------------------------
% Si DocSetup n'a pas défini \DocDisplayBoxPictures, on met un défaut.
% Si une ancienne macro \displayboxpictures existe, on l'utilise comme source.
\makeatletter
\@ifundefined{DocDisplayBoxPictures}{%
	\@ifundefined{displayboxpictures}{%
		% ni DocDisplayBoxPictures ni displayboxpictures -> défaut
		\providecommand{\DocDisplayBoxPictures}{../_includes/boxes_images}%
	}{%
		% ancienne macro displayboxpictures définie -> on la réutilise
		\let\DocDisplayBoxPictures\displayboxpictures%
	}%
}{}%
\makeatother

%----------------------------------------------------------------------
% Packages nécessaires (tolérant si déjà chargés ailleurs)
%----------------------------------------------------------------------
\RequirePackage{graphicx}
\RequirePackage{marginnote}
\RequirePackage{tcolorbox}
\RequirePackage{tikz}
\RequirePackage{verbatim}
\RequirePackage{etoolbox}

% Librairies tcolorbox utiles
\tcbuselibrary{skins,breakable}

%======================================================================
% Macro utilitaire pour déclarer les boîtes standard
%======================================================================
% \CFPTDeclareBox{<NomEnv>}{<prefixCouleur>}{<colTitreTexte>}{<épaisseurCadre>}{<icone>}
%
% Paramètres :
%  - <NomEnv>        : nom de l'environnement tcolorbox à créer
%                      (boxInfo, boxSuccess, boxWarning, …).
%  - <prefixCouleur> : préfixe des couleurs définies dans inc_colors :
%                      <prefix>Background, <prefix>Color, <prefix>TitleBackground.
%                      Ex. "info", "success", "warning", "danger", "teacher",
%                      "download", "aide", "objectif", "validation", "default".
%  - <colTitreTexte> : couleur du texte du titre (ex. white, black).
%  - <épaisseurCadre>: épaisseur du cadre (ex. 0.5pt, 1.0pt).
%  - <icone>         : nom de fichier icône (ex. box-info.png) dans
%                      \DocDisplayBoxPictures. Chaîne vide ("") = aucune icône.
%
% Comportement spécifique :
%  - Les boîtes sont "breakable" (coupure possible sur plusieurs pages).
%  - L'icône est dessinée en overlay en haut à droite :
%      * une fois pour une boîte non cassée (overlay unbroken),
%      * une fois sur le premier segment pour une boîte cassée (overlay first).
%  - La marge droite (clé 'right') dépend de la présence d'icône et de titre :
%      * pas d'icône (#5 vide)              -> right = 0.5cm
%      * icône + titre présent              -> right = 0.5cm
%      * icône mais PAS de titre            -> right = 1.0cm
%    -> le texte ne passe pas sous l'icône lorsqu'il n'y a pas de titre.
%======================================================================

\makeatletter
\newcommand{\CFPTDeclareBox}[5]{%
	\newtcolorbox{#1}[1][]{%
		colback=#2Background,
		colframe=#2Color,
		colbacktitle=#2TitleBackground,
		coltitle=#3,
		boxrule=#4,
		width=\linewidth,
		enhanced,
		breakable,
		%--------------------------------------------
		% Marge à droite :
		%  - pas d'icône (#5 vide)           -> 0mm
		%  - icône + titre présent          -> 0mm
		%  - icône mais PAS de titre        -> 1.2cm
		%--------------------------------------------
		right=\ifx&#5&%
					0.5cm%
				\else
					\iftcb@hasTitle%
						0.5cm%
					\else
						1.0cm%
					\fi
				\fi,
		title={##1},
		fonttitle=\bfseries,
		toptitle=2mm,
		bottomtitle=2mm,
		% cacher l'entête si pas de titre
		title code={\if\relax\detokenize{##1}\relax\else\tcbtitle\par\fi},
		% Icône : un seul dessin (boîte non cassée ou 1er segment)
		overlay unbroken={%
			\ifx&#5&%
			\else
			\node[anchor=north east] at (frame.north east)
			{\includegraphics[height=0.70cm]{\DocDisplayBoxPictures/#5}};%
			\fi
		},
		overlay first={%
			\ifx&#5&%
			\else
			\node[anchor=north east] at (frame.north east)
			{\includegraphics[height=0.70cm]{\DocDisplayBoxPictures/#5}};%
			\fi
		},
	}%
}
\makeatother

%======================================================================
% Déclaration des boîtes standard (couleurs centralisées dans inc_colors)
%======================================================================

% boîte par défaut grise, sans icône
\CFPTDeclareBox{boxDefault}{default}{white}{0.5pt}{}

% info
\CFPTDeclareBox{boxInfo}{info}{white}{0.5pt}{box-info.png}

% success
\CFPTDeclareBox{boxSuccess}{success}{white}{0.5pt}{box-success.png}

% warning (titre en noir pour lisibilité sur fond jaune)
\CFPTDeclareBox{boxWarning}{warning}{black}{0.5pt}{box-warning.png}

% danger
\CFPTDeclareBox{boxDanger}{danger}{white}{0.5pt}{box-danger.png}

% teacher
\CFPTDeclareBox{boxTeacher}{teacher}{white}{0.5pt}{box-teacher.png}

% download (cadre plus épais)
\CFPTDeclareBox{boxDownload}{download}{white}{1.0pt}{box-download.png}

% aide
\CFPTDeclareBox{boxHelp}{aide}{white}{0.5pt}{box-aide.png}

% Objectif
\CFPTDeclareBox{boxObjectif}{objectif}{black}{0.5pt}{box-objectif.png}

% Validation
\CFPTDeclareBox{boxValidation}{validation}{black}{0.5pt}{box-validation.png}


% ======================================
% TOUS LES LANGAGES
% ======================================

% Web


% ==============================
% CONFIGURATION GLOBALE LISTINGS
% ==============================

\usepackage{listings}
\usepackage{xcolor}
\usepackage{inconsolata}

% ------------------------------------------------
% Palette moderne (inspirée GitHub)
% ------------------------------------------------
\definecolor{codeBg}{HTML}{F6F8FA}
\definecolor{codeText}{HTML}{24292F}
\definecolor{codeFrame}{HTML}{D0D7DE}
\definecolor{codeComment}{HTML}{6A737D}
\definecolor{codeString}{HTML}{032F62}
\definecolor{codeNumber}{HTML}{005CC5}

% Couleurs par langage
\definecolor{phpColor}{HTML}{6F42C1}
\definecolor{jsColor}{HTML}{D97706}
\definecolor{htmlColor}{HTML}{C92C2C}
\definecolor{cssColor}{HTML}{005CC5}
\definecolor{sqlColor}{HTML}{22863A}
\definecolor{javaColor}{HTML}{B07219}
\definecolor{csharpColor}{HTML}{178600}

% ------------------------------------------------
% Style général commun à TOUS les langages
% ------------------------------------------------
\lstdefinestyle{styleCode}{
  backgroundcolor=\color{codeBg},
  basicstyle=\ttfamily\small\color{codeText},
  frame=single,
  rulecolor=\color{codeFrame},
  numbers=left,
  numberstyle=\tiny\color{codeComment},
  stepnumber=1,
  numbersep=8pt,
  tabsize=2,
  showspaces=false,
  showstringspaces=false,
  showtabs=false,
  breaklines=true,
  breakatwhitespace=true,
  captionpos=b,
  aboveskip=1.2em,
  belowskip=1.2em,
  commentstyle=\color{codeComment}\itshape,
  stringstyle=\color{codeString},
  identifierstyle=\color{codeText},
  keywordstyle=\bfseries,
  columns=fullflexible,
  keepspaces=true,
  inputencoding=utf8,
  extendedchars=true,
  literate={
    {é}{{\'e}}1
    {à}{{\`a}}1
    {è}{{\`e}}1
    {ù}{{\`u}}1
    {â}{{\^a}}1
    {ê}{{\^e}}1
    {î}{{\^i}}1
    {ô}{{\^o}}1
    {û}{{\^u}}1
    {ë}{{\"e}}1
    {ï}{{\"i}}1
    {ö}{{\"o}}1
    {ü}{{\"u}}1
    {✅}{{\ding{51}}}1
    {❌}{{\ding{53}}}1
    {🔐}{{}}1
    {📱}{{}}1
    {🚀}{{}}1
    {📦}{{}}1
    {⚠️}{{}}1
    {●}{{}}1
    % {<}{{\<}}1
    % {>}{{\>}}1
    {É}{{\'E}}1
    {|}{{\textbar}}1
    {œ}{{\oe}}1
    {ç}{{\c{c}}}1
  }
}

% ======================================
% HTML
% ======================================

\lstdefinelanguage{HTMLcustom}{
  language=HTML,
  keywordstyle=\color{htmlColor}\bfseries,
  sensitive=false
}

\lstnewenvironment{htmlcode}[1][]{
  \lstset{
    style=styleCode,
    language=HTMLcustom,
    #1
  }
}{}


% ==============================
% CONFIGURATION GLOBALE LISTINGS
% ==============================

\usepackage{listings}
\usepackage{xcolor}
\usepackage{inconsolata}

% ------------------------------------------------
% Palette moderne (inspirée GitHub)
% ------------------------------------------------
\definecolor{codeBg}{HTML}{F6F8FA}
\definecolor{codeText}{HTML}{24292F}
\definecolor{codeFrame}{HTML}{D0D7DE}
\definecolor{codeComment}{HTML}{6A737D}
\definecolor{codeString}{HTML}{032F62}
\definecolor{codeNumber}{HTML}{005CC5}

% Couleurs par langage
\definecolor{phpColor}{HTML}{6F42C1}
\definecolor{jsColor}{HTML}{D97706}
\definecolor{htmlColor}{HTML}{C92C2C}
\definecolor{cssColor}{HTML}{005CC5}
\definecolor{sqlColor}{HTML}{22863A}
\definecolor{javaColor}{HTML}{B07219}
\definecolor{csharpColor}{HTML}{178600}

% ------------------------------------------------
% Style général commun à TOUS les langages
% ------------------------------------------------
\lstdefinestyle{styleCode}{
  backgroundcolor=\color{codeBg},
  basicstyle=\ttfamily\small\color{codeText},
  frame=single,
  rulecolor=\color{codeFrame},
  numbers=left,
  numberstyle=\tiny\color{codeComment},
  stepnumber=1,
  numbersep=8pt,
  tabsize=2,
  showspaces=false,
  showstringspaces=false,
  showtabs=false,
  breaklines=true,
  breakatwhitespace=true,
  captionpos=b,
  aboveskip=1.2em,
  belowskip=1.2em,
  commentstyle=\color{codeComment}\itshape,
  stringstyle=\color{codeString},
  identifierstyle=\color{codeText},
  keywordstyle=\bfseries,
  columns=fullflexible,
  keepspaces=true,
  inputencoding=utf8,
  extendedchars=true,
  literate={
    {é}{{\'e}}1
    {à}{{\`a}}1
    {è}{{\`e}}1
    {ù}{{\`u}}1
    {â}{{\^a}}1
    {ê}{{\^e}}1
    {î}{{\^i}}1
    {ô}{{\^o}}1
    {û}{{\^u}}1
    {ë}{{\"e}}1
    {ï}{{\"i}}1
    {ö}{{\"o}}1
    {ü}{{\"u}}1
    {✅}{{\ding{51}}}1
    {❌}{{\ding{53}}}1
    {🔐}{{}}1
    {📱}{{}}1
    {🚀}{{}}1
    {📦}{{}}1
    {⚠️}{{}}1
    {●}{{}}1
    % {<}{{\<}}1
    % {>}{{\>}}1
    {É}{{\'E}}1
    {|}{{\textbar}}1
    {œ}{{\oe}}1
    {ç}{{\c{c}}}1
  }
}

% ======================================
% CSS
% ======================================

\lstdefinelanguage{CSScustom}{
  keywords={
    align-items, background, border, border-radius,
    box-shadow, color, display, flex, font-size,
    font-weight, height, justify-content, margin,
    padding, position, transition, transform,
    width, max-width, min-width
  },
  keywordstyle=\color{cssColor}\bfseries,
  sensitive=true,
  morecomment=[l]{/*},
  morecomment=[s]{/*}{*/},
  morestring=[b]",
  alsoletter={-},
}

\lstnewenvironment{csscode}[1][]{
  \lstset{
    style=styleCode,
    language=CSScustom,
    #1
  }
}{}


% ==============================
% CONFIGURATION GLOBALE LISTINGS
% ==============================

\usepackage{listings}
\usepackage{xcolor}
\usepackage{inconsolata}

% ------------------------------------------------
% Palette moderne (inspirée GitHub)
% ------------------------------------------------
\definecolor{codeBg}{HTML}{F6F8FA}
\definecolor{codeText}{HTML}{24292F}
\definecolor{codeFrame}{HTML}{D0D7DE}
\definecolor{codeComment}{HTML}{6A737D}
\definecolor{codeString}{HTML}{032F62}
\definecolor{codeNumber}{HTML}{005CC5}

% Couleurs par langage
\definecolor{phpColor}{HTML}{6F42C1}
\definecolor{jsColor}{HTML}{D97706}
\definecolor{htmlColor}{HTML}{C92C2C}
\definecolor{cssColor}{HTML}{005CC5}
\definecolor{sqlColor}{HTML}{22863A}
\definecolor{javaColor}{HTML}{B07219}
\definecolor{csharpColor}{HTML}{178600}

% ------------------------------------------------
% Style général commun à TOUS les langages
% ------------------------------------------------
\lstdefinestyle{styleCode}{
  backgroundcolor=\color{codeBg},
  basicstyle=\ttfamily\small\color{codeText},
  frame=single,
  rulecolor=\color{codeFrame},
  numbers=left,
  numberstyle=\tiny\color{codeComment},
  stepnumber=1,
  numbersep=8pt,
  tabsize=2,
  showspaces=false,
  showstringspaces=false,
  showtabs=false,
  breaklines=true,
  breakatwhitespace=true,
  captionpos=b,
  aboveskip=1.2em,
  belowskip=1.2em,
  commentstyle=\color{codeComment}\itshape,
  stringstyle=\color{codeString},
  identifierstyle=\color{codeText},
  keywordstyle=\bfseries,
  columns=fullflexible,
  keepspaces=true,
  inputencoding=utf8,
  extendedchars=true,
  literate={
    {é}{{\'e}}1
    {à}{{\`a}}1
    {è}{{\`e}}1
    {ù}{{\`u}}1
    {â}{{\^a}}1
    {ê}{{\^e}}1
    {î}{{\^i}}1
    {ô}{{\^o}}1
    {û}{{\^u}}1
    {ë}{{\"e}}1
    {ï}{{\"i}}1
    {ö}{{\"o}}1
    {ü}{{\"u}}1
    {✅}{{\ding{51}}}1
    {❌}{{\ding{53}}}1
    {🔐}{{}}1
    {📱}{{}}1
    {🚀}{{}}1
    {📦}{{}}1
    {⚠️}{{}}1
    {●}{{}}1
    % {<}{{\<}}1
    % {>}{{\>}}1
    {É}{{\'E}}1
    {|}{{\textbar}}1
    {œ}{{\oe}}1
    {ç}{{\c{c}}}1
  }
}

% ======================================
% JAVASCRIPT
% ======================================

\lstdefinelanguage{JScustom}{
  keywords={
    break, case, catch, class, const, continue, debugger,
    default, delete, do, else, export, extends, finally,
    for, function, if, import, in, instanceof, let,
    new, return, super, switch, this, throw, try,
    typeof, var, void, while, with, yield,
    async, await
  },
  keywordstyle=\color{jsColor}\bfseries,
  sensitive=true,
  comment=[l]{//},
  morecomment=[s]{/*}{*/}
}

\lstnewenvironment{jscode}[1][]{
  \lstset{
    style=styleCode,
    language=JScustom,
    #1
  }
}{}


% ==============================
% CONFIGURATION GLOBALE LISTINGS
% ==============================

\usepackage{listings}
\usepackage{xcolor}
\usepackage{inconsolata}

% ------------------------------------------------
% Palette moderne (inspirée GitHub)
% ------------------------------------------------
\definecolor{codeBg}{HTML}{F6F8FA}
\definecolor{codeText}{HTML}{24292F}
\definecolor{codeFrame}{HTML}{D0D7DE}
\definecolor{codeComment}{HTML}{6A737D}
\definecolor{codeString}{HTML}{032F62}
\definecolor{codeNumber}{HTML}{005CC5}

% Couleurs par langage
\definecolor{phpColor}{HTML}{6F42C1}
\definecolor{jsColor}{HTML}{D97706}
\definecolor{htmlColor}{HTML}{C92C2C}
\definecolor{cssColor}{HTML}{005CC5}
\definecolor{sqlColor}{HTML}{22863A}
\definecolor{javaColor}{HTML}{B07219}
\definecolor{csharpColor}{HTML}{178600}

% ------------------------------------------------
% Style général commun à TOUS les langages
% ------------------------------------------------
\lstdefinestyle{styleCode}{
  backgroundcolor=\color{codeBg},
  basicstyle=\ttfamily\small\color{codeText},
  frame=single,
  rulecolor=\color{codeFrame},
  numbers=left,
  numberstyle=\tiny\color{codeComment},
  stepnumber=1,
  numbersep=8pt,
  tabsize=2,
  showspaces=false,
  showstringspaces=false,
  showtabs=false,
  breaklines=true,
  breakatwhitespace=true,
  captionpos=b,
  aboveskip=1.2em,
  belowskip=1.2em,
  commentstyle=\color{codeComment}\itshape,
  stringstyle=\color{codeString},
  identifierstyle=\color{codeText},
  keywordstyle=\bfseries,
  columns=fullflexible,
  keepspaces=true,
  inputencoding=utf8,
  extendedchars=true,
  literate={
    {é}{{\'e}}1
    {à}{{\`a}}1
    {è}{{\`e}}1
    {ù}{{\`u}}1
    {â}{{\^a}}1
    {ê}{{\^e}}1
    {î}{{\^i}}1
    {ô}{{\^o}}1
    {û}{{\^u}}1
    {ë}{{\"e}}1
    {ï}{{\"i}}1
    {ö}{{\"o}}1
    {ü}{{\"u}}1
    {✅}{{\ding{51}}}1
    {❌}{{\ding{53}}}1
    {🔐}{{}}1
    {📱}{{}}1
    {🚀}{{}}1
    {📦}{{}}1
    {⚠️}{{}}1
    {●}{{}}1
    % {<}{{\<}}1
    % {>}{{\>}}1
    {É}{{\'E}}1
    {|}{{\textbar}}1
    {œ}{{\oe}}1
    {ç}{{\c{c}}}1
  }
}

% ======================================
% PHP
% ======================================

\lstdefinelanguage{PHPcustom}{
  language=PHP,
  morekeywords={
    self, parent, static,
    public, private, protected,
    use, function
  },
  keywordstyle=\color{phpColor}\bfseries,
  sensitive=true
}

\lstnewenvironment{phpcode}[1][]{
  \lstset{
    style=styleCode,
    language=PHPcustom,
    #1
  }
}{}

% Autre


% ==============================
% CONFIGURATION GLOBALE LISTINGS
% ==============================

\usepackage{listings}
\usepackage{xcolor}
\usepackage{inconsolata}

% ------------------------------------------------
% Palette moderne (inspirée GitHub)
% ------------------------------------------------
\definecolor{codeBg}{HTML}{F6F8FA}
\definecolor{codeText}{HTML}{24292F}
\definecolor{codeFrame}{HTML}{D0D7DE}
\definecolor{codeComment}{HTML}{6A737D}
\definecolor{codeString}{HTML}{032F62}
\definecolor{codeNumber}{HTML}{005CC5}

% Couleurs par langage
\definecolor{phpColor}{HTML}{6F42C1}
\definecolor{jsColor}{HTML}{D97706}
\definecolor{htmlColor}{HTML}{C92C2C}
\definecolor{cssColor}{HTML}{005CC5}
\definecolor{sqlColor}{HTML}{22863A}
\definecolor{javaColor}{HTML}{B07219}
\definecolor{csharpColor}{HTML}{178600}

% ------------------------------------------------
% Style général commun à TOUS les langages
% ------------------------------------------------
\lstdefinestyle{styleCode}{
  backgroundcolor=\color{codeBg},
  basicstyle=\ttfamily\small\color{codeText},
  frame=single,
  rulecolor=\color{codeFrame},
  numbers=left,
  numberstyle=\tiny\color{codeComment},
  stepnumber=1,
  numbersep=8pt,
  tabsize=2,
  showspaces=false,
  showstringspaces=false,
  showtabs=false,
  breaklines=true,
  breakatwhitespace=true,
  captionpos=b,
  aboveskip=1.2em,
  belowskip=1.2em,
  commentstyle=\color{codeComment}\itshape,
  stringstyle=\color{codeString},
  identifierstyle=\color{codeText},
  keywordstyle=\bfseries,
  columns=fullflexible,
  keepspaces=true,
  inputencoding=utf8,
  extendedchars=true,
  literate={
    {é}{{\'e}}1
    {à}{{\`a}}1
    {è}{{\`e}}1
    {ù}{{\`u}}1
    {â}{{\^a}}1
    {ê}{{\^e}}1
    {î}{{\^i}}1
    {ô}{{\^o}}1
    {û}{{\^u}}1
    {ë}{{\"e}}1
    {ï}{{\"i}}1
    {ö}{{\"o}}1
    {ü}{{\"u}}1
    {✅}{{\ding{51}}}1
    {❌}{{\ding{53}}}1
    {🔐}{{}}1
    {📱}{{}}1
    {🚀}{{}}1
    {📦}{{}}1
    {⚠️}{{}}1
    {●}{{}}1
    % {<}{{\<}}1
    % {>}{{\>}}1
    {É}{{\'E}}1
    {|}{{\textbar}}1
    {œ}{{\oe}}1
    {ç}{{\c{c}}}1
  }
}

% ======================================
% SQL
% ======================================

\lstdefinelanguage{SQLcustom}{
  language=SQL,
  keywordstyle=\color{sqlColor}\bfseries,
  sensitive=false
}

\lstnewenvironment{sqlcode}[1][]{
  \lstset{
    style=styleCode,
    language=SQLcustom,
    #1
  }
}{}


% ==============================
% CONFIGURATION GLOBALE LISTINGS
% ==============================

\usepackage{listings}
\usepackage{xcolor}
\usepackage{inconsolata}

% ------------------------------------------------
% Palette moderne (inspirée GitHub)
% ------------------------------------------------
\definecolor{codeBg}{HTML}{F6F8FA}
\definecolor{codeText}{HTML}{24292F}
\definecolor{codeFrame}{HTML}{D0D7DE}
\definecolor{codeComment}{HTML}{6A737D}
\definecolor{codeString}{HTML}{032F62}
\definecolor{codeNumber}{HTML}{005CC5}

% Couleurs par langage
\definecolor{phpColor}{HTML}{6F42C1}
\definecolor{jsColor}{HTML}{D97706}
\definecolor{htmlColor}{HTML}{C92C2C}
\definecolor{cssColor}{HTML}{005CC5}
\definecolor{sqlColor}{HTML}{22863A}
\definecolor{javaColor}{HTML}{B07219}
\definecolor{csharpColor}{HTML}{178600}

% ------------------------------------------------
% Style général commun à TOUS les langages
% ------------------------------------------------
\lstdefinestyle{styleCode}{
  backgroundcolor=\color{codeBg},
  basicstyle=\ttfamily\small\color{codeText},
  frame=single,
  rulecolor=\color{codeFrame},
  numbers=left,
  numberstyle=\tiny\color{codeComment},
  stepnumber=1,
  numbersep=8pt,
  tabsize=2,
  showspaces=false,
  showstringspaces=false,
  showtabs=false,
  breaklines=true,
  breakatwhitespace=true,
  captionpos=b,
  aboveskip=1.2em,
  belowskip=1.2em,
  commentstyle=\color{codeComment}\itshape,
  stringstyle=\color{codeString},
  identifierstyle=\color{codeText},
  keywordstyle=\bfseries,
  columns=fullflexible,
  keepspaces=true,
  inputencoding=utf8,
  extendedchars=true,
  literate={
    {é}{{\'e}}1
    {à}{{\`a}}1
    {è}{{\`e}}1
    {ù}{{\`u}}1
    {â}{{\^a}}1
    {ê}{{\^e}}1
    {î}{{\^i}}1
    {ô}{{\^o}}1
    {û}{{\^u}}1
    {ë}{{\"e}}1
    {ï}{{\"i}}1
    {ö}{{\"o}}1
    {ü}{{\"u}}1
    {✅}{{\ding{51}}}1
    {❌}{{\ding{53}}}1
    {🔐}{{}}1
    {📱}{{}}1
    {🚀}{{}}1
    {📦}{{}}1
    {⚠️}{{}}1
    {●}{{}}1
    % {<}{{\<}}1
    % {>}{{\>}}1
    {É}{{\'E}}1
    {|}{{\textbar}}1
    {œ}{{\oe}}1
    {ç}{{\c{c}}}1
  }
}

% ======================================
% C# / CSHARP
% ======================================

\lstdefinelanguage{CSharpcustom}{
  language=[Sharp]C,
  keywordstyle=\color{csharpColor}\bfseries,
  sensitive=true
}

\lstnewenvironment{csharpcode}[1][]{
  \lstset{
    style=styleCode,
    language=CSharpcustom,
    #1
  }
}{}


% ==============================
% CONFIGURATION GLOBALE LISTINGS
% ==============================

\usepackage{listings}
\usepackage{xcolor}
\usepackage{inconsolata}

% ------------------------------------------------
% Palette moderne (inspirée GitHub)
% ------------------------------------------------
\definecolor{codeBg}{HTML}{F6F8FA}
\definecolor{codeText}{HTML}{24292F}
\definecolor{codeFrame}{HTML}{D0D7DE}
\definecolor{codeComment}{HTML}{6A737D}
\definecolor{codeString}{HTML}{032F62}
\definecolor{codeNumber}{HTML}{005CC5}

% Couleurs par langage
\definecolor{phpColor}{HTML}{6F42C1}
\definecolor{jsColor}{HTML}{D97706}
\definecolor{htmlColor}{HTML}{C92C2C}
\definecolor{cssColor}{HTML}{005CC5}
\definecolor{sqlColor}{HTML}{22863A}
\definecolor{javaColor}{HTML}{B07219}
\definecolor{csharpColor}{HTML}{178600}

% ------------------------------------------------
% Style général commun à TOUS les langages
% ------------------------------------------------
\lstdefinestyle{styleCode}{
  backgroundcolor=\color{codeBg},
  basicstyle=\ttfamily\small\color{codeText},
  frame=single,
  rulecolor=\color{codeFrame},
  numbers=left,
  numberstyle=\tiny\color{codeComment},
  stepnumber=1,
  numbersep=8pt,
  tabsize=2,
  showspaces=false,
  showstringspaces=false,
  showtabs=false,
  breaklines=true,
  breakatwhitespace=true,
  captionpos=b,
  aboveskip=1.2em,
  belowskip=1.2em,
  commentstyle=\color{codeComment}\itshape,
  stringstyle=\color{codeString},
  identifierstyle=\color{codeText},
  keywordstyle=\bfseries,
  columns=fullflexible,
  keepspaces=true,
  inputencoding=utf8,
  extendedchars=true,
  literate={
    {é}{{\'e}}1
    {à}{{\`a}}1
    {è}{{\`e}}1
    {ù}{{\`u}}1
    {â}{{\^a}}1
    {ê}{{\^e}}1
    {î}{{\^i}}1
    {ô}{{\^o}}1
    {û}{{\^u}}1
    {ë}{{\"e}}1
    {ï}{{\"i}}1
    {ö}{{\"o}}1
    {ü}{{\"u}}1
    {✅}{{\ding{51}}}1
    {❌}{{\ding{53}}}1
    {🔐}{{}}1
    {📱}{{}}1
    {🚀}{{}}1
    {📦}{{}}1
    {⚠️}{{}}1
    {●}{{}}1
    % {<}{{\<}}1
    % {>}{{\>}}1
    {É}{{\'E}}1
    {|}{{\textbar}}1
    {œ}{{\oe}}1
    {ç}{{\c{c}}}1
  }
}

% ======================================
% JAVA
% ======================================

\lstdefinelanguage{Javacustom}{
  language=Java,
  keywordstyle=\color{javaColor}\bfseries,
  sensitive=true
}

\lstnewenvironment{javacode}[1][]{
  \lstset{
    style=styleCode,
    language=Javacustom,
    #1
  }
}{}

% --- packages ---
\usepackage{forest}
\usepackage{fontawesome5}

% --- encadré esthétique (optionnel) ---
% \newtcolorbox{treebox}{
% 	colback=white, colframe=black!10, boxrule=0.4pt,
% 	borderline west={2pt}{0pt}{blue!35},
% 	arc=2mm, left=8pt, right=8pt, top=8pt, bottom=8pt
% }

% --- palette ---
\definecolor{pill}{HTML}{EEF2FF}
\definecolor{pillborder}{HTML}{C7D2FE}
\definecolor{filepill}{HTML}{F3F4F6}
\definecolor{fileborder}{HTML}{E5E7EB}
\definecolor{ink}{HTML}{0F172A}

% --- interrupteur + marge dépendante ---
\newif\ificons
\newdimen\iconpad
% réglage global simple : \SetIcons{true} ou \SetIcons{false}
\newcommand{\SetIcons}[1]{%
	\csname icons#1\endcsname
	\ificons \iconpad=1.1em\relax \else \iconpad=0.3em\relax \fi
}
% init : sans icônes
\SetIcons{false}

% --- base commune forest ---
\forestset{
	base layout/.style={
		for tree={
			font=\sffamily\small\color{ink},
			grow'=0,
			parent anchor=east, child anchor=west,
			edge path'={(!u.parent anchor) -- +(7pt,0) |- (.child anchor)},
			edge={draw=black!25, line width=.5pt},
			l sep=12pt,
			s sep=7pt,
			anchor=west, calign=first,
			text height=1.2ex, text depth=.2ex,
			inner xsep=\iconpad,  % marge commune (dépend de \SetIcons)
		},
		pill/.style={draw, rounded corners=4pt, minimum height=1.4em},
		folder/.style={pill, fill=pill,     draw=pillborder},
		file/.style  ={pill, fill=filepill, draw=fileborder},
	},
	compact/.style={for tree/.append style={l sep=8pt, s sep=5pt}},
}

% --- icônes : fallback si rien n'est indiqué ---
\newcommand{\DefaultIcon}{file} % icône par défaut

% case d'icône : réserve la place ; affiche l'icône si \iconstrue
\newcommand{\IconBox}[1]{%
	\makebox[\iconpad][c]{\ificons \faIcon{#1}\fi}%
}
% version “peut-être vide” : si argument vide => DefaultIcon
\newcommand{\IconBoxMaybe}[1]{%
	\if\relax\detokenize{#1}\relax
	\IconBox{\DefaultIcon}%
	\else
	\IconBox{#1}%
	\fi
}

% --- macros de label (renommées) ---
\newcommand{\froot}[1]{\IconBox{folder-open}#1}
\newcommand{\fdir}[1]{\IconBox{folder}#1}
\newcommand{\ffile}[1]{\IconBox{file}#1}
\newcommand{\fcode}[1]{\IconBox{file-code}#1}
\newcommand{\fimage}[1]{\IconBox{image}#1} % ou "file-image"

% (optionnel) générique : \withicon[<nom-ou-vide>]{Texte}
\newcommand{\withicon}[2][]{\IconBoxMaybe{#1}#2}

% Config en-têtes
\renewcommand{\sectionmark}[1]{\markright{#1}}
\renewcommand{\chaptermark}[1]{\markboth{\thechapter\ -\ #1}{}}
\newcommand{\currentSection}{\rightmark}
\newcommand{\currentChapiter}{\leftmark}


% En-têtes et pieds de page
\newcommand{\leftHead}{\myTitle}
\newcommand{\centerHead}{\myLesson}
\newcommand{\rightHead}{\currentChapiter}

\newcommand{\leftFoot}{\myAuthor \\ {\small \myEmail}}
\newcommand{\centerFoot}{\mySchool\\{\textcolor{cfptGray}{\small Généré le \today}}}
\newcommand{\rightFoot}{\thepage~/~\pageref{LastPage}}


% Definition des en-têtes et pieds de page
\fancyhead[L]{\leftHead}
\fancyhead[C]{\centerHead}
\fancyhead[R]{\rightHead}

\fancyfoot[L]{\leftFoot}
\fancyfoot[C]{\centerFoot}
\fancyfoot[R]{\rightFoot}

% Redéfinition des en-têtes et pieds pour la page de chapitre
\fancypagestyle{plain}{
	\fancyhead[L]{}
	\fancyhead[C]{}
	\fancyhead[R]{}
}


% Lignes de séparation dans les en-têtes/pieds
\renewcommand{\headrulewidth}{0.4pt}
\renewcommand{\footrulewidth}{0.4pt}

% Commandes


% -------------------------
% Commandes personnalisées
% -------------------------

% Base
\newcommand{\newLine}{\vspace{0.3cm}}
\newcommand{\newMediumLine}{\vspace{0.175cm}}
\newcommand{\newSmallLine}{\vspace{0.05cm}}
\newcommand{\bdd}{base de données}
\newcommand{\tpi}{Travail Pratique Individuel}

% Legendes de tests
\newcommand{\legendItem}[2]{#1\hspace{0.5em}#2}
\newcommand{\gap}{\space\space\space\space\space\space}

% Etats de tests
\newcommand{\good}{\textcolor{green!60!black}{\ding{51}}}
\newcommand{\notGood}{\textcolor{red!70!black}{\ding{55}}}
\newcommand{\notImpl}{\textcolor{gray}{--}}
